\QHeading{第1题　脱式计算。}
\par\noindent


\CalcBlock[方法一]{1}{\begin{aligned}[t]&\phantom{={}}\color{black}460\div[(34+12)\times2]\\&\color{answerred}=460\div[46\times2]\\&\color{answerred}=460\div92\\&\color{answerred}=5\end{aligned}}
\par
\MethodBlock{方法二：除以积，改成连续除}{连续除以两个数，等于除以这两个数的积。}{\begin{aligned}[t]&\phantom{={}}460\div [(34+12)\times 2]\\
&=460\div (34+12)\div 2\\
&=460\div 46\div 2\\
&=10\div 2\\
&=5\end{aligned}}


\CalcBlock{2}{\begin{aligned}[t]&\phantom{={}}\color{black}52\times[288\div(130-106)]\\&\color{answerred}=52\times[288\div24]\\&\color{answerred}=52\times12\\&\color{answerred}=624\end{aligned}}
\CalcBlock[方法一]{3}{\begin{aligned}[t]&\phantom{={}}\color{black}350\div[7\times(45-40)]\\&\color{answerred}=350\div[7\times5]\\&\color{answerred}=350\div35\\&\color{answerred}=10\end{aligned}}
\MethodBlock{方法二：除以积，改成连续除}{连续除以两个数，等于除以这两个数的积。}{\begin{aligned}[t]&\phantom{={}}350\div [7\times (45-40)]\\
&=350\div 7\div (45-40)\\
&=350\div 7\div 5\\
&=50\div 5\\
&=10\end{aligned}}

\CalcBlock{4}{\begin{aligned}[t]&\phantom{={}}\color{black}140\div[(56-28)\div4]\\&\color{answerred}=140\div[28\div4]\\&\color{answerred}=140\div7\\&\color{answerred}=20\end{aligned}}
\CalcBlock[简便解法（补充）]{5}{\begin{aligned}[t]&\phantom{={}}\color{black}[70-(82-60)]\times 25\\&\color{answerred}=[70-22]\times 25\\&\color{answerred}=48\times 25\\&\color{answerred}=(12\times 4)\times 25\\&\color{answerred}=12\times (4\times 25)\\&\color{answerred}=12\times 100\\&\color{answerred}=1200\end{aligned}}
\MethodBlock{按运算顺序计算（基础练习）}{本单元还要练习运算顺序；与上面的简便解法对照。}{\begin{aligned}[t]&\phantom{={}}[70-(82-60)]\times25\\&=[70-22]\times25\\&=48\times25\\&=1200\end{aligned}}

\MethodBlock{方法二：乘法分配律}{分别求出两部分的积，再合起来。}{\begin{aligned}[t]&\phantom{={}}[70-(82-60)]\times 25\\
&=70\times 25-(82-60)\times 25\\
&=70\times 25-22\times 25\\
&=1750-550\\
&=1200\end{aligned}}

\CalcBlock[方法一]{6}{\begin{aligned}[t]&\phantom{={}}\color{black}42\times[169-(78+35)]\\&\color{answerred}=42\times[169-113]\\&\color{answerred}=42\times56\\&\color{answerred}=2352\end{aligned}}
\MethodBlock{方法二：乘法分配律}{把括号中的两部分分别相乘，再相加或相减。}{\begin{aligned}[t]&\phantom{={}}42\times [169-(78+35)]\\
&=42\times 169-42\times (78+35)\\
&=42\times 169-42\times 113\\
&=7098-4746\\
&=2352\end{aligned}}

\par\noindent

\QHeading{第2题　比一比，算一算。}

\CalcBlock{1}{\begin{aligned}[t]&\phantom{={}}\color{black}960\div40-10\div2\\&\color{answerred}=24-5\\&\color{answerred}=19\end{aligned}}
\par

\CalcBlock[方法一]{2}{\begin{aligned}[t]&\phantom{={}}\color{black}(960\div40-10)\div2\\&\color{answerred}=(24-10)\div2\\&\color{answerred}=14\div2\\&\color{answerred}=7\end{aligned}}
\MethodBlock{方法二：先分别相除再合并}{}{\begin{aligned}[t]&\phantom{={}}(960\div 40-10)\div 2\\
&=960\div 40\div 2-10\div 2\\
&=24\div 2-5\\
&=12-5\\
&=7\end{aligned}}

\CalcBlock{3}{\begin{aligned}[t]&\phantom{={}}\color{black}960\div[(40-10)\div2]\\&\color{answerred}=960\div[30\div2]\\&\color{answerred}=960\div15\\&\color{answerred}=64\end{aligned}}
\par\noindent

\QHeading{第3题　解决问题。}

\Problem{①王老师去书店买书，每套故事书35元，每套科技书85元，王老师买了15套故事书和15套科技书，买科技书比买故事书多花多少钱？}
\par


\MethodTitle{方法一}
\Answer{\Formula{\begin{aligned}[t]&\phantom{={}}(85-35)\times15\\&=50\times15\\&=750\text{（元）}\end{aligned}}\par 答：买科技书比买故事书多花750元。}

\MethodTitle{方法二：分别求两类书的总价}
\Answer{先分别计算科技书和故事书的费用，再相减。}
\Answer{\Formula{\begin{aligned}[t]&\phantom{={}}85\times 15-35\times 15\\
&=1275-525\\
&=750\text{（元）}\end{aligned}}}
\Answer{答：买科技书比买故事书多花750元。}
\Problem{②光明小学举行演讲比赛。三年级有40人参赛，四年级的参赛人数比三年级多25人，五年级的参赛人数比四年级的2倍少15人。五年级有多少人参赛？}

\MethodTitle{方法一}
\Answer{\Formula{\begin{aligned}[t]&\phantom{={}}(40+25)\times2-15\\&=65\times2-15\\&=130-15\\&=115\text{（人）}\end{aligned}}\par 答：五年级有115人参赛。}

\MethodTitle{方法二：把四年级人数分成两部分}
\Answer{四年级是40人再加25人；它的2倍就是两部分分别乘2，再减15。}
\Answer{\Formula{\begin{aligned}[t]&\phantom{={}}40\times 2+25\times 2-15\\
&=80+50-15\\
&=130-15\\
&=115\text{（人）}\end{aligned}}}
\Answer{答：五年级有115人参赛。}
